% Options for packages loaded elsewhere
% Options for packages loaded elsewhere
\PassOptionsToPackage{unicode}{hyperref}
\PassOptionsToPackage{hyphens}{url}
\PassOptionsToPackage{dvipsnames,svgnames,x11names}{xcolor}
%
\documentclass[
  spanish,
  letterpaper,
  DIV=11,
  numbers=noendperiod]{scrartcl}
\usepackage{xcolor}
\usepackage{amsmath,amssymb}
\setcounter{secnumdepth}{-\maxdimen} % remove section numbering
\usepackage{iftex}
\ifPDFTeX
  \usepackage[T1]{fontenc}
  \usepackage[utf8]{inputenc}
  \usepackage{textcomp} % provide euro and other symbols
\else % if luatex or xetex
  \usepackage{unicode-math} % this also loads fontspec
  \defaultfontfeatures{Scale=MatchLowercase}
  \defaultfontfeatures[\rmfamily]{Ligatures=TeX,Scale=1}
\fi
\usepackage{lmodern}
\ifPDFTeX\else
  % xetex/luatex font selection
\fi
% Use upquote if available, for straight quotes in verbatim environments
\IfFileExists{upquote.sty}{\usepackage{upquote}}{}
\IfFileExists{microtype.sty}{% use microtype if available
  \usepackage[]{microtype}
  \UseMicrotypeSet[protrusion]{basicmath} % disable protrusion for tt fonts
}{}
\makeatletter
\@ifundefined{KOMAClassName}{% if non-KOMA class
  \IfFileExists{parskip.sty}{%
    \usepackage{parskip}
  }{% else
    \setlength{\parindent}{0pt}
    \setlength{\parskip}{6pt plus 2pt minus 1pt}}
}{% if KOMA class
  \KOMAoptions{parskip=half}}
\makeatother
% Make \paragraph and \subparagraph free-standing
\makeatletter
\ifx\paragraph\undefined\else
  \let\oldparagraph\paragraph
  \renewcommand{\paragraph}{
    \@ifstar
      \xxxParagraphStar
      \xxxParagraphNoStar
  }
  \newcommand{\xxxParagraphStar}[1]{\oldparagraph*{#1}\mbox{}}
  \newcommand{\xxxParagraphNoStar}[1]{\oldparagraph{#1}\mbox{}}
\fi
\ifx\subparagraph\undefined\else
  \let\oldsubparagraph\subparagraph
  \renewcommand{\subparagraph}{
    \@ifstar
      \xxxSubParagraphStar
      \xxxSubParagraphNoStar
  }
  \newcommand{\xxxSubParagraphStar}[1]{\oldsubparagraph*{#1}\mbox{}}
  \newcommand{\xxxSubParagraphNoStar}[1]{\oldsubparagraph{#1}\mbox{}}
\fi
\makeatother

\usepackage{color}
\usepackage{fancyvrb}
\newcommand{\VerbBar}{|}
\newcommand{\VERB}{\Verb[commandchars=\\\{\}]}
\DefineVerbatimEnvironment{Highlighting}{Verbatim}{commandchars=\\\{\}}
% Add ',fontsize=\small' for more characters per line
\usepackage{framed}
\definecolor{shadecolor}{RGB}{241,243,245}
\newenvironment{Shaded}{\begin{snugshade}}{\end{snugshade}}
\newcommand{\AlertTok}[1]{\textcolor[rgb]{0.68,0.00,0.00}{#1}}
\newcommand{\AnnotationTok}[1]{\textcolor[rgb]{0.37,0.37,0.37}{#1}}
\newcommand{\AttributeTok}[1]{\textcolor[rgb]{0.40,0.45,0.13}{#1}}
\newcommand{\BaseNTok}[1]{\textcolor[rgb]{0.68,0.00,0.00}{#1}}
\newcommand{\BuiltInTok}[1]{\textcolor[rgb]{0.00,0.23,0.31}{#1}}
\newcommand{\CharTok}[1]{\textcolor[rgb]{0.13,0.47,0.30}{#1}}
\newcommand{\CommentTok}[1]{\textcolor[rgb]{0.37,0.37,0.37}{#1}}
\newcommand{\CommentVarTok}[1]{\textcolor[rgb]{0.37,0.37,0.37}{\textit{#1}}}
\newcommand{\ConstantTok}[1]{\textcolor[rgb]{0.56,0.35,0.01}{#1}}
\newcommand{\ControlFlowTok}[1]{\textcolor[rgb]{0.00,0.23,0.31}{\textbf{#1}}}
\newcommand{\DataTypeTok}[1]{\textcolor[rgb]{0.68,0.00,0.00}{#1}}
\newcommand{\DecValTok}[1]{\textcolor[rgb]{0.68,0.00,0.00}{#1}}
\newcommand{\DocumentationTok}[1]{\textcolor[rgb]{0.37,0.37,0.37}{\textit{#1}}}
\newcommand{\ErrorTok}[1]{\textcolor[rgb]{0.68,0.00,0.00}{#1}}
\newcommand{\ExtensionTok}[1]{\textcolor[rgb]{0.00,0.23,0.31}{#1}}
\newcommand{\FloatTok}[1]{\textcolor[rgb]{0.68,0.00,0.00}{#1}}
\newcommand{\FunctionTok}[1]{\textcolor[rgb]{0.28,0.35,0.67}{#1}}
\newcommand{\ImportTok}[1]{\textcolor[rgb]{0.00,0.46,0.62}{#1}}
\newcommand{\InformationTok}[1]{\textcolor[rgb]{0.37,0.37,0.37}{#1}}
\newcommand{\KeywordTok}[1]{\textcolor[rgb]{0.00,0.23,0.31}{\textbf{#1}}}
\newcommand{\NormalTok}[1]{\textcolor[rgb]{0.00,0.23,0.31}{#1}}
\newcommand{\OperatorTok}[1]{\textcolor[rgb]{0.37,0.37,0.37}{#1}}
\newcommand{\OtherTok}[1]{\textcolor[rgb]{0.00,0.23,0.31}{#1}}
\newcommand{\PreprocessorTok}[1]{\textcolor[rgb]{0.68,0.00,0.00}{#1}}
\newcommand{\RegionMarkerTok}[1]{\textcolor[rgb]{0.00,0.23,0.31}{#1}}
\newcommand{\SpecialCharTok}[1]{\textcolor[rgb]{0.37,0.37,0.37}{#1}}
\newcommand{\SpecialStringTok}[1]{\textcolor[rgb]{0.13,0.47,0.30}{#1}}
\newcommand{\StringTok}[1]{\textcolor[rgb]{0.13,0.47,0.30}{#1}}
\newcommand{\VariableTok}[1]{\textcolor[rgb]{0.07,0.07,0.07}{#1}}
\newcommand{\VerbatimStringTok}[1]{\textcolor[rgb]{0.13,0.47,0.30}{#1}}
\newcommand{\WarningTok}[1]{\textcolor[rgb]{0.37,0.37,0.37}{\textit{#1}}}

\usepackage{longtable,booktabs,array}
\newcounter{none} % for unnumbered tables
\usepackage{calc} % for calculating minipage widths
% Correct order of tables after \paragraph or \subparagraph
\usepackage{etoolbox}
\makeatletter
\patchcmd\longtable{\par}{\if@noskipsec\mbox{}\fi\par}{}{}
\makeatother
% Allow footnotes in longtable head/foot
\IfFileExists{footnotehyper.sty}{\usepackage{footnotehyper}}{\usepackage{footnote}}
\makesavenoteenv{longtable}
\usepackage{graphicx}
\makeatletter
\newsavebox\pandoc@box
\newcommand*\pandocbounded[1]{% scales image to fit in text height/width
  \sbox\pandoc@box{#1}%
  \Gscale@div\@tempa{\textheight}{\dimexpr\ht\pandoc@box+\dp\pandoc@box\relax}%
  \Gscale@div\@tempb{\linewidth}{\wd\pandoc@box}%
  \ifdim\@tempb\p@<\@tempa\p@\let\@tempa\@tempb\fi% select the smaller of both
  \ifdim\@tempa\p@<\p@\scalebox{\@tempa}{\usebox\pandoc@box}%
  \else\usebox{\pandoc@box}%
  \fi%
}
% Set default figure placement to htbp
\def\fps@figure{htbp}
\makeatother



\ifLuaTeX
\usepackage[bidi=basic,shorthands=off]{babel}
\else
\usepackage[bidi=default,shorthands=off]{babel}
\fi
\ifLuaTeX
  \usepackage{selnolig} % disable illegal ligatures
\fi


\setlength{\emergencystretch}{3em} % prevent overfull lines

\providecommand{\tightlist}{%
  \setlength{\itemsep}{0pt}\setlength{\parskip}{0pt}}



 


\usepackage{booktabs}
\usepackage{longtable}
\usepackage{array}
\usepackage{multirow}
\usepackage{wrapfig}
\usepackage{float}
\usepackage{colortbl}
\usepackage{pdflscape}
\usepackage{tabularx}
\usepackage{xltabular}
\usepackage{threeparttable}
\usepackage{threeparttablex}
\usepackage[normalem]{ulem}
\usepackage{makecell}
\usepackage{xcolor}
\KOMAoption{captions}{tableheading}
\makeatletter
\@ifpackageloaded{caption}{}{\usepackage{caption}}
\AtBeginDocument{%
\ifdefined\contentsname
  \renewcommand*\contentsname{Tabla de contenidos}
\else
  \newcommand\contentsname{Tabla de contenidos}
\fi
\ifdefined\listfigurename
  \renewcommand*\listfigurename{Listado de Figuras}
\else
  \newcommand\listfigurename{Listado de Figuras}
\fi
\ifdefined\listtablename
  \renewcommand*\listtablename{Listado de Tablas}
\else
  \newcommand\listtablename{Listado de Tablas}
\fi
\ifdefined\figurename
  \renewcommand*\figurename{Figura}
\else
  \newcommand\figurename{Figura}
\fi
\ifdefined\tablename
  \renewcommand*\tablename{Tabla}
\else
  \newcommand\tablename{Tabla}
\fi
}
\@ifpackageloaded{float}{}{\usepackage{float}}
\floatstyle{ruled}
\@ifundefined{c@chapter}{\newfloat{codelisting}{h}{lop}}{\newfloat{codelisting}{h}{lop}[chapter]}
\floatname{codelisting}{Listado}
\newcommand*\listoflistings{\listof{codelisting}{Listado de Listados}}
\makeatother
\makeatletter
\makeatother
\makeatletter
\@ifpackageloaded{caption}{}{\usepackage{caption}}
\@ifpackageloaded{subcaption}{}{\usepackage{subcaption}}
\makeatother
\usepackage{bookmark}
\IfFileExists{xurl.sty}{\usepackage{xurl}}{} % add URL line breaks if available
\urlstyle{same}
% fallback for those not using the hyperref driver hyperxmp:
\makeatletter
\@ifundefined{xmpquote}{\newcommand{\xmpquote}[1]{#1}}{}
\makeatother
\hypersetup{
  pdftitle={ACP: Análisis de Componentes Principales},
  pdflang={es},
  colorlinks=true,
  linkcolor={blue},
  filecolor={Maroon},
  citecolor={Blue},
  urlcolor={Blue},
  pdfcreator={LaTeX via pandoc}}


\title{ACP: Análisis de Componentes Principales}
\author{}
\date{}
\begin{document}
\maketitle


\textbf{AUTORES}

Jared Mijail Ramírez Escalante · María Paula Silva Capera

PDF

\subsection{Introducción}\label{introducciuxf3n}

A continuación se presenta el análisis de componentes principales
implementado sobre las variables cuantitativas del conjunto de datos
PetFinder.

\subsubsection{Carga de librerías y
datos}\label{carga-de-libreruxedas-y-datos}

\begin{Shaded}
\begin{Highlighting}[]
\FunctionTok{library}\NormalTok{(here)}
\FunctionTok{library}\NormalTok{(ggplot2)}
\FunctionTok{library}\NormalTok{(plotly)}
\FunctionTok{library}\NormalTok{(ggrepel)}
\FunctionTok{library}\NormalTok{(factoextra)}
\FunctionTok{library}\NormalTok{(FactoMineR)}
\FunctionTok{library}\NormalTok{(naniar)}
\FunctionTok{library}\NormalTok{(knitr)}
\FunctionTok{library}\NormalTok{(kableExtra)}
\FunctionTok{library}\NormalTok{(stringr)}

\CommentTok{\# Cargamos la tabla de datos enriquecida}
\NormalTok{petfinder\_enriquecido\_final }\OtherTok{\textless{}{-}} \FunctionTok{read.csv}\NormalTok{(}
  \FunctionTok{file.path}\NormalTok{(}\FunctionTok{here}\NormalTok{(), }\StringTok{"data"}\NormalTok{, }\StringTok{"petfinder\_enriquecido\_final.csv"}\NormalTok{)}
\NormalTok{)}

\NormalTok{n }\OtherTok{\textless{}{-}} \FunctionTok{nrow}\NormalTok{(petfinder\_enriquecido\_final)}
\end{Highlighting}
\end{Shaded}

\subsection{Selección de variables activas y
suplementarias}\label{selecciuxf3n-de-variables-activas-y-suplementarias}

\subsubsection{Variables activas}\label{variables-activas}

\begin{Shaded}
\begin{Highlighting}[]
\NormalTok{vars\_activas }\OtherTok{\textless{}{-}} \FunctionTok{c}\NormalTok{(}
  \StringTok{\textquotesingle{}PhotoAmt\textquotesingle{}}\NormalTok{,}
  \StringTok{\textquotesingle{}sentiment\_score\textquotesingle{}}\NormalTok{,}
  \StringTok{\textquotesingle{}sentiment\_magnitude\textquotesingle{}}\NormalTok{,}
  \StringTok{\textquotesingle{}total\_labels\textquotesingle{}}\NormalTok{,}
  \StringTok{\textquotesingle{}total\_colors\textquotesingle{}}\NormalTok{,}
  \StringTok{\textquotesingle{}desc\_char\_len\textquotesingle{}}\NormalTok{,}
  \StringTok{\textquotesingle{}desc\_word\_count\textquotesingle{}}
\NormalTok{)}
\end{Highlighting}
\end{Shaded}

Las variables elegidas para ser partícipes de la construcción de los
ejes del plano factorial nuevo fueron \emph{PhotoAmt},
\emph{sentiment\_score}, \emph{sentiment\_magnitude},
\emph{total\_labels}, \emph{total\_colors}, \emph{desc\_char\_len} y
\emph{desc\_word\_count}. Todas las anteriores tienen en común que son
métricas asociadas al anuncio elaborado por los rescatistas sobre las
mascotas que ofrecen en adopción mediante la plataforma PetFinder. La
razón por la cual se decide emplear esta agrupación de variables, es
porque las variables \emph{PhotoAmt}, \emph{total\_labels} y
\emph{total\_colors} se refieren a las características de las fotos que
subieron los rescatistas al anuncio, mientras que las variables
\emph{sentiment\_score}, \emph{sentiment\_magnitude},
\emph{desc\_char\_len} y \emph{desc\_word\_count} están relacionadas con
el texto que escribieron. Las 7 variables se relacionan con la
elaboración y contenido del anuncio, características que reflejan la
cantidad y el tipo de información que los rescatistas incorporaron en la
elaboración del anuncio.

\subsubsection{Variables suplementarias}\label{variables-suplementarias}

\begin{Shaded}
\begin{Highlighting}[]
\NormalTok{vars\_sup }\OtherTok{\textless{}{-}} \FunctionTok{c}\NormalTok{(}\StringTok{\textquotesingle{}Age\textquotesingle{}}\NormalTok{, }\StringTok{\textquotesingle{}Quantity\textquotesingle{}}\NormalTok{, }\StringTok{\textquotesingle{}Fee\textquotesingle{}}\NormalTok{)}
\NormalTok{cuantitativas }\OtherTok{\textless{}{-}} \FunctionTok{c}\NormalTok{(vars\_activas, vars\_sup)}

\CommentTok{\# Identificadores de las mascotas}
\NormalTok{ids }\OtherTok{\textless{}{-}}\NormalTok{ petfinder\_enriquecido\_final}\SpecialCharTok{$}\NormalTok{PetID}
\FunctionTok{rownames}\NormalTok{(petfinder\_enriquecido\_final) }\OtherTok{\textless{}{-}}\NormalTok{ ids}

\CommentTok{\# Conjunto de datos con las variables cuantitativas}
\NormalTok{petfinder\_cuantis }\OtherTok{\textless{}{-}}\NormalTok{ petfinder\_enriquecido\_final[, cuantitativas]}
\end{Highlighting}
\end{Shaded}

Las variables suplementarias \emph{Age}, \emph{Quantity} y \emph{Fee}
representan características externas a la construcción del anuncio, son
propias del animal o de las condiciones de adopción y se proyectan
posteriormente para evaluar su relación con la estructura identificada.

\subsection{Tratamiento de los valores
faltantes}\label{tratamiento-de-los-valores-faltantes}

Como se evidenció en el diagnóstico univariado de la entrega 01, hay
variables que presentan valores faltantes. Para poder ejecutar el
análisis de componentes principales, es importante darles un tratamiento
adecuado para que la matriz numérica esté completa. Dentro de nuestras
variables cuantitativas.

\begin{Shaded}
\begin{Highlighting}[]
\CommentTok{\# Número total de valores faltantes}
\FunctionTok{sum}\NormalTok{(}\FunctionTok{is.na}\NormalTok{(petfinder\_enriquecido\_final))}
\end{Highlighting}
\end{Shaded}

\begin{verbatim}
[1] 3382
\end{verbatim}

\begin{Shaded}
\begin{Highlighting}[]
\CommentTok{\# Variables con valores faltantes}
\NormalTok{vars\_na }\OtherTok{\textless{}{-}} \FunctionTok{names}\NormalTok{(}
  \FunctionTok{which}\NormalTok{(}\FunctionTok{colSums}\NormalTok{(}\FunctionTok{is.na}\NormalTok{(petfinder\_cuantis)) }\SpecialCharTok{\textgreater{}} \DecValTok{0}\NormalTok{)}
\NormalTok{)}

\NormalTok{na\_df }\OtherTok{\textless{}{-}} \FunctionTok{data.frame}\NormalTok{(}
  \AttributeTok{variable =} \FunctionTok{names}\NormalTok{(petfinder\_cuantis),}
  \AttributeTok{n\_na     =} \FunctionTok{colSums}\NormalTok{(}\FunctionTok{is.na}\NormalTok{(petfinder\_cuantis)),}
  \AttributeTok{row.names =} \ConstantTok{NULL}
\NormalTok{)}
\NormalTok{knitr}\SpecialCharTok{::}\FunctionTok{kable}\NormalTok{(na\_df, }\AttributeTok{col.names =} \FunctionTok{c}\NormalTok{(}\StringTok{"Variable"}\NormalTok{, }\StringTok{"Valores faltantes"}\NormalTok{))}
\end{Highlighting}
\end{Shaded}

{\def\LTcaptype{none} % do not increment counter
\begin{longtable}[]{@{}lr@{}}
\toprule\noalign{}
Variable & Valores faltantes \\
\midrule\noalign{}
\endhead
\bottomrule\noalign{}
\endlastfoot
PhotoAmt & 0 \\
sentiment\_score & 551 \\
sentiment\_magnitude & 551 \\
total\_labels & 341 \\
total\_colors & 341 \\
desc\_char\_len & 0 \\
desc\_word\_count & 0 \\
Age & 0 \\
Quantity & 0 \\
Fee & 0 \\
\end{longtable}
}

Se identificó que \emph{total\_labels}, \emph{total\_colors},
\emph{sentiment\_score} y \emph{sentiment\_magnitude} presentan valores
faltantes. Todas las anteriores son variables derivadas del
enriquecimiento de imágen y de texto realizado con la API de Google
Cloud. Para indagar sobre la ocurrencia de estos valores faltantes y
posibles patrones, primero se realiza la prueba estadística de Little
para poner a prueba la hipótesis nula H0: los datos faltantes son MCAR
(Missing Completely At Random) versus la alternativa Ha: los datos
faltantes no son MCAR.

\begin{Shaded}
\begin{Highlighting}[]
\CommentTok{\# Prueba de Little}
\CommentTok{\# H0: los datos son MCAR}
\CommentTok{\# H1: los datos no son MCAR}

\FunctionTok{mcar\_test}\NormalTok{(petfinder\_cuantis)}
\end{Highlighting}
\end{Shaded}

\begin{verbatim}
# A tibble: 1 x 4
  statistic    df p.value missing.patterns
      <dbl> <dbl>   <dbl>            <int>
1     1264.    22       0                4
\end{verbatim}

Se obtuvo un valor p de 0 entonces se rechaza que los datos sean MCAR.
Esto sugiere que sí esxiste una aparición sistemática de los valores
faltantes, por lo tanto es pertinente realizar una verificación de cómo
son estos posibles patrones. Para ello es útil visualizar cuántos
valores faltantes existen en cada una de las variables respecto a las
otras y respecto al total, y cuántos valores faltantes comparten entre
ellas

\begin{Shaded}
\begin{Highlighting}[]
\CommentTok{\# Patrones de ocurrencia}
\FunctionTok{gg\_miss\_upset}\NormalTok{(petfinder\_cuantis[, vars\_na])}
\end{Highlighting}
\end{Shaded}

\pandocbounded{\includegraphics[keepaspectratio]{PCA_files/figure-pdf/unnamed-chunk-5-1.pdf}}

Se observa que el par de variables \emph{total\_labels} y
\emph{total\_colors} comparten 341 individuos con valores faltantes
mientras que el par \emph{sentiment\_score} y
\emph{sentiment\_magnitude} comparten 551. Hay 331 individuos que
carecen de observaciones para las variables \emph{total\_labels} y
\emph{total\_colors}, 541 que no tienen valores en
\emph{sentiment\_score} y \emph{sentiment\_magnitude} y 10 individuos
que no tienen un valor asignado para ninguna de estas 4 variables. Esto
es consistente con el hecho de que las variables \emph{total\_labels} y
\emph{total\_colors} dependen del enriquecimiento de imagen, mientras
que las variables \emph{sentiment\_socre} y \emph{sentiment\_magnitude}
dependen del enriquecimiento de texto. Lo anterior sugiere que hubo
individuos cuyas imágenes y/o texto no fueron correctamente procesados
por la API de Google Cloud. Para verificarlo se visualizó la tabla de
los individuos que presentan valores faltantes para estos pares de
variables.




\end{document}
